\subsection{インダストリー4.0とソフトウェア工学}
インダストリー4.0は、製造業を大きく変える概念である。特に、AIを活用した個別化された製造、デジタル化、標準システムとの統合が重要になる。これにより、費用、品質、効率に便益が期待される。
インダストリー4.0では、ソフトウェアが中心的な構成要素になる。多くの機器が相互接続され、データを送り、命令やデータを受け取り、さらに動作する。したがって、堅牢で知的なシステムを作るには、ソフトウェアを工学的に作ることが不可欠である。
継続的ソフトウェア工学は、継続的な計画、アーキテクチャ設計、開発、統合、配備、レビュー・改訂を支える。製造が継続的に変化するなら、ソフトウェア開発も継続的に学習し、更新する必要がある。
\par\smallskip\noindent\refstepcounter{table}\label{tab:ch18-14}表 \thetable: インダストリー4.0とソフトウェア工学の整理\par\smallskip
表 \thetable は、インダストリー4.0とソフトウェア工学の整理を比較・確認しやすい形で整理したものである。\par\smallskip
\begin{longtable}{>{\raggedright\arraybackslash}p{0.28\linewidth}>{\raggedright\arraybackslash}p{0.64\linewidth}}
\toprule
技術 & ソフトウェア工学での意味 \\
\midrule
\endfirsthead
\toprule
技術 & ソフトウェア工学での意味 \\
\midrule
\endhead
IoT & 機器からデータを集め、機器へ命令を返すソフトウェアが必要になる。 \\
ビッグデータ分析 & 大量データを処理し、予測や最適化に使う。 \\
AI・機械学習 & 予測、検査、自動判断、個別化された製造に使う。 \\
サイバーセキュリティ & 接続機器が増えるほど攻撃面も増えるため、設計から保護が必要である。 \\
クラウドコンピューティング & 計算資源、保存、監視、分析を柔軟に提供する。 \\
複数プラットフォーム向けアプリ & 工場、モバイル、クラウド、組込み機器をまたいだ利用者体験を支える。 \\
量子コンピューティング & 複雑な計算をより高速かつ費用対効果高く実行する可能性がある。 \\
\bottomrule
\end{longtable}
今後のソフトウェアは、自己学習的で先回りする性質を強める可能性がある。利用者が何を望むかを予測し、状況に応じて振る舞いを変えるためには、測定、実験、モデル、標準、根本原因分析のすべてが重要になる。
\par\smallskip
\begin{center}
\centering
\IfFileExists{assets/generated_figures/ch18/fig-ch18-16.png}{\includegraphics[width=0.96\linewidth]{assets/generated_figures/ch18/fig-ch18-16.png}}{\fbox{\parbox[c][48mm][c]{0.9\linewidth}{\centering 画像生成待ち\\インダストリー4.0とソフトウェア工学の関係図}}}
\par\smallskip
\refstepcounter{figure}\label{fig:ch18-16}図 \thefigure: インダストリー4.0とソフトウェア工学の関係図
\end{center}
図 \thefigure は、インダストリー4.0とソフトウェア工学の関係図を視覚的に整理したものである。
\par\smallskip
